\subsection{The comparison and gluing propositions}

\begin{lemma}[Artinian comparison on all test schemes]\label{A:comparison}
\status{Proved}
Let $S$ be locally noetherian, let $G,F$ be Zariski and fpqc sheaves
commuting with filtered ring colimits and complete noetherian local
adic limits, and let $v:G\to F$ be a natural transformation. If
$v(A)$ is bijective for every Artinian local $S$-algebra $A$, then
$v$ is an isomorphism on all $S$-schemes.
\end{lemma}
\begin{proof}
We first prove injectivity on all rings. Over a noetherian ring, two
sections of $G$ with the same image become equal modulo every power
of every maximal ideal by the Artinian hypothesis. Adic compatibility
makes them equal over the complete local rings; fpqc injectivity makes
them equal over the local rings. Filtered-colimit compatibility then
spreads equality to neighbourhoods, and the Zariski sheaf property
gives equality globally. Every algebra over a noetherian affine base
is a filtered union of finite-type subalgebras; an element pair and
its equality in $F$ descend to one such subalgebra. This proves
injectivity for arbitrary rings.

For surjectivity over a noetherian local ring $A$, an element of
$F(A)$ has unique, compatible preimages in $G(A/\mathfrak m^n)$.
They give an element of $G(\widehat A)$, with the prescribed image.
Its two pullbacks to $G(\widehat A\otimes_A\widehat A)$ agree by
the already established injectivity on \emph{arbitrary} rings.
Fpqc descent gives an element of $G(A)$. Its image is the original
element, checked after the faithfully flat completion. As in the
injectivity argument, finite-type approximation and Zariski gluing
extend surjectivity first to noetherian rings and then to all rings
and schemes. Every comparison is the natural pullback map; no
noetherianity of $\widehat A\otimes_A\widehat A$ is used.
\end{proof}

For $G=h_X$ with $X$ locally of finite presentation, the hypotheses on
$G$ hold. In particular, a compatible system of maps from the
Artinian quotients of a complete local ring lands in any fixed affine
neighbourhood of its common closed-point image. Taking limits of the
coordinate maps gives the required map to $X$. This observation
explains why this adic assertion does not require $X$ to be proper.
Lemma~\ref{A:comparison} supplies the missing general-argument
injectivity in cote 131, scans 27--29, and rechecks the substantive
claim of the supplied R1.

\begin{proposition}[The formal gluing proposition]\label{A:formal-gluing}
\status{Proved under the displayed overlap hypothesis}
Let $A$ be complete noetherian local, let $S=\Spec A$, let $s$ be
its closed point, and put $U=S-\{s\}$. Let $F$ satisfy conditions
1--4. Suppose we are given a locally finite-type scheme $X_U$
representing $F_U$, finite local $A$-algebras $B_i$, and compatible
elements $\xi_i\in F(B_i)$, with the following properties:
\begin{enumerate}
\item the morphism
$\coprod_i(\Spec B_i)_U\to X_U$ induced by $\xi_i$ is an open
immersion;
\item the maps $\Spec B_i\to F$ together represent $F$ on all
Artinian local $A$-algebras whose structural closed point maps to $s$.
\end{enumerate}
Then $F$ is represented by the scheme obtained by gluing
$\coprod_i\Spec B_i$ to $X_U$ along that open immersion.
\end{proposition}
\begin{proof}
We glue the two schemes along their specified open subschemes, using
ordinary scheme gluing (Stacks, Lemma 26.14.2, Tag 01JC)
\cite{Stacks01JC}. The result $X$ is locally of finite type over $S$.
Condition 1 glues the given sections to $\xi\in F(X)$ and hence
defines $h_X\to F$. For an Artinian local $A$-algebra $C$, either
its point maps to $s$, in which case (2) applies, or its structural
map factors through $U$, in which case the representing $X_U$
applies. Thus $h_X(C)\to F(C)$ is bijective for every such $C$.
Lemma~\ref{A:comparison} proves $h_X\simeq F$ on all schemes.
The hypothesis (1) is exactly the condition singled out in
Proposition 2 on scans 31--33; this proof does not assert that
Artinian pro-modularity alone establishes (1).
\end{proof}

\begin{lemma}[A finiteness fact for local isomorphisms]\label{A:localiso-qc}
\status{Proved}
A separated local isomorphism $P\to T$ with finite fibres is
quasi-compact if $T$ is noetherian. Here a local isomorphism means
that every point of $P$ has an open neighbourhood mapped
isomorphically onto an open of $T$.
\end{lemma}
\begin{proof}
Suppose first that $T$ is irreducible, with generic point $\eta$.
Every such neighbourhood in $P$ contains a unique point above
$\eta$. Two neighbourhoods with the same point above $\eta$
give two sections over the intersection of their images. Their
equalizer is open because the morphism is a local isomorphism,
and closed because it is separated. It contains the generic point
of the irreducible intersection, so it is the entire intersection.
Thus all neighbourhoods belonging to one generic point glue to
a single open of $P$ isomorphic to an open of $T$. These opens
are disjoint and cover $P$. There are finitely many of them since
the fibre over $\eta$ is finite. Every open of a noetherian space
is quasi-compact, proving the claim in this case.

For general $T$, choose a finite partition into irreducible locally
closed subschemes. Such a partition exists by noetherian induction
on closed subsets. The preceding argument makes the inverse image
of each stratum quasi-compact. Their finitely many underlying
subspaces cover $P$, so every open cover of $P$ has a finite
subcover. Hence $P\to T$ is quasi-compact: the same argument applies
over each affine open of the noetherian $T$.
\end{proof}

\begin{theorem}[Proposition 3, with its gluing argument completed]\label{A:local-to-global}
\status{Proved under precisely stated local representability and constructibility}
Let $S$ be locally noetherian and let $F$ be a Zariski sheaf commuting
with filtered ring colimits. Suppose:
\begin{enumerate}
\item for every $s\in S$, the restriction $F_{S_s}$, where
$S_s=\Spec\OO_{S,s}$, is represented on all $S_s$-schemes by a
scheme $X_s$ locally of finite type;
\item for every locally finite-type $S$-scheme $T$ and
$\eta\in F(T)$, the locus $M(T,\eta)$ of local modularity is
locally constructible.
\end{enumerate}
Then $F$ is represented by an $S$-scheme locally of finite type.
The charts called ``quasi-modular'' at the end of cote 131 are
open immersions $h_V\hookrightarrow F$.
\end{theorem}
\begin{proof}
We give all three passages: openness, representable local charts,
and removal of the extra identifications.

\emph{Openness.} The locus $M(T,\eta)$ is stable under
generization. Indeed, if $t'$ generalizes $t$ and $s$ is the image
of $t$, both points lie over $S_s$. In the representing $X_s$,
local modularity at $t$ says exactly that its local-ring map is an
isomorphism. Localizing that isomorphism at the prime corresponding
to $t'$ gives local modularity at $t'$. A locally constructible
generization-stable subset of a locally noetherian scheme is open.
Thus every $M(T,\eta)$ is open.

\emph{Spreading a chart.} Fix $s\in S$ and an affine open
$V_s\subset X_s$ of finite type over $S_s$. Shrink $S$ to an affine
neighbourhood of $s$. Finite presentation spreads $V_s$ to an affine
finite-type $S$-scheme $V$, and condition 3 spreads its universal
element to $\xi\in F(V)$. The locus $M(V,\xi)$ contains all of
$V_{S_s}$. Its closed complement is a finite-type scheme with
empty base change to $S_s$, so it is empty after shrinking $S$
around $s$. To check the last assertion explicitly, cover that
complement by finitely many affine algebras. Each becomes the zero
ring after localization at $s$; finitely many denominators outside
the prime of $s$ therefore kill all of them. We may hence assume
every point of $V$ is locally modular.

\emph{Representability of the chart map.} Let $T$ be any
$S$-scheme and let $\eta\in F(T)$. Put
$P=h_V\times_F h_T$. For each $t\in T$, local representability
over $S_{f(t)}$ represents $P_{T_t}$ by a local isomorphism over
$T_t=\Spec\OO_{T,t}$. This holds because the maps on local rings
of $V$ to the corresponding $X_{f(t)}$ are isomorphisms; finite
presentation turns these local isomorphisms of rings into
isomorphisms on neighbourhoods.

Every point of the closed fibre of $P_{T_t}$ is $\kappa(t)$-rational
and extends to a section over $T_t$: its local-isomorphism
neighbourhood has open image containing the closed point of the
local scheme, hence has image all of $T_t$. Since $P$ commutes
with filtered ring colimits, that section spreads to an open
neighbourhood of $t$ in $T$. Two such sections agree on an open
subset: agreement at a point implies agreement over its local
scheme by the local-isomorphism property just proved, and filtered
colimits spread that agreement to a neighbourhood. We therefore
glue the resulting copies of opens of $T$ along their open
equalizers to form a scheme $P_0$ locally isomorphic to $T$.

This gluing represents the full functor $P$, not just its sections
on opens of $T$. In fact, for a section over an arbitrary
$T'\to T$, its value at $t'\in T'$ lies on a unique branch of the
closed fibre of $P_{T_t}$ after extension to $\kappa(t')$.
That fibre is a disjoint union of copies of $\Spec\kappa(t)$,
so the branch is already defined over $\kappa(t)$ and extends
to one of our sections near $t$. Local uniqueness at the same
point of $V$ identifies this section with the given one over
$T'_{t'}$, and the equality spreads. Zariski locality then
identifies $h_{P_0}$ with $P$ on every $T'$. Thus $h_V\to F$
is representable by local isomorphisms.

\emph{Removing the extra identifications.} We apply this to
$R=h_V\times_F h_V$. Its first projection is a local isomorphism.
The morphism $R\to V\times_S V$ is a monomorphism, hence separated;
as $V\to S$ is affine, $R\to V$ is separated. Its fibres are
finite. To see this at a point of $V$ lying over $b\in S$, use
the representing $X_b$: the fibre is a fibre of the local
isomorphism $V_{S_b}\to X_b$. Its source is noetherian, so that
morphism is quasi-compact, locally of finite type, and has
zero-dimensional fibres. Hence those fibres are finite.
Lemma~\ref{A:localiso-qc} makes $R$ quasi-compact. It is consequently
of finite type over the noetherian affine neighbourhood of $s$.

The diagonal section $V\to R$ is both open and closed, since
$R\to V$ is separated and a local isomorphism. Its complement
$R-\Delta_V$ is of finite type and has empty base change to
$S_s$, because $V_s\to F_{S_s}$ is an open immersion. The same
finite-denominator argument shrinks $S$ so that this complement
is empty. Hence $h_V\to F$ is a monomorphism as well as a
representable local isomorphism. It is therefore an open immersion.

Taking all the above charts covers $F$. Indeed, for every
$\eta\in F(T)$ and $t\in T$, the corresponding point of $X_{f(t)}$
lies in some affine $V_{f(t)}$. The constructed open chart has a
lift near $t$ by the preceding local-isomorphism argument. The
charts and their represented open overlaps glue as schemes, and
the Zariski sheaf property identifies the glued functor with $F$,
as in Stacks, Lemma 26.15.4, Tag 01JJ \cite{Stacks01JJ}.
Local finite type follows on these charts.

Finally, the boxed insertion on scan 41 explicitly describes
the desired map $h_{X^{(x)}}\to F$ as an open immersion. We have
inspected that scan. This is the precise meaning needed by its
subsequent gluing sentence, rather than a synonym for an
unproved constructibility assertion.
\end{proof}

The hypotheses in Theorem~\ref{A:local-to-global} are the actual
two inputs of the manuscript's Proposition 3 on scans 33--41.
They are not asserted to follow from the eight axioms by the
proof just given. In particular the restriction of a modularity
locus to a closed stratum need not equal the modularity locus
after that closed base change. That distinction is the remaining
issue in the proposed noetherian-induction lead L1.
